\begin{abstract}
Direct Preference Optimization (DPO) has emerged as an efficient alternative to PPO-based RLHF, yet its application to knowledge-intensive generation tasks reveals a fundamental limitation: standard preference signals---whether from human annotators or LLM judges---exhibit a systematic \textit{verbosity bias} that rewards linguistic fluency over logical correctness. We show empirically that this reward signal blindspot leaves a large logical alignment gap: SFT-trained models achieve NLI entailment scores of only 0.05--0.22, despite producing fluent, confident-sounding text. We propose \textbf{RLearner-LLM}, a framework that resolves this through \textbf{Hybrid-DPO}: an automated preference construction pipeline that fuses a DeBERTa-v3 NLI entailment signal with a verifier LLM score, eliminating the need for human annotation while overcoming the ``alignment tax'' that degrades fluency under single-signal optimization. Evaluated across five academic domains (Biology, Medicine, Law) with two architectures (LLaMA-2-13B, Qwen3-8B), RLearner-LLM achieves up to 6$\times$ NLI improvement and consistent answer coverage gains over SFT baselines. Our Qwen3-8B model achieves a 95\% win rate against both the SFT baseline and GPT-4o-mini in blind pairwise evaluation, demonstrating that a locally-hosted 8B model aligned with our hybrid reward can match proprietary frontier models on educational explanation quality.
\end{abstract}
